% =====================================================================
%  Chapter 30 opener -- "Chapter at a Glance", three reading paths + the
%  ladder staircase, THE running example (the 6-vertex labelled tree T on
%  {1,2,3,4,5,6} with edges {1-5, 2-5, 5-3, 3-4, 4-6} and Prufer code
%  [5,5,3,4]), the Cayley / Matrix-Tree / Expectation taxonomy, the two
%  traps (labelled vs unlabelled, and the "code is sequence" trap that
%  misses the bijection), the who-owns-what comparison against Ch.28
%  (unlabelled shapes), Ch.13 (tree DP), Ch.19 (segment tree), Ch.12
%  (depth/path), Ch.17 (DSU), Ch.31 (MST -- the chapter AFTER next), and
%  Ch.29 (Trees Born from Graphs -- the PREVIOUS chapter), a constraint
%  -> technique table, and notation.
%  Exports fig:c32-tree, fig:c32-code, fig:c32-taxonomy.
%  \term the NEW objects "Prufer code", "labelled tree", "spanning-tree
%  census", "Laplacian", "cofactor", "matrix-tree", "linearity of expec-
%  tation on trees" only -- degree/depth/DSU are indexed by earlier chapters.
% =====================================================================

\begin{chaptertheorybox}{Chapter at a Glance}
Every earlier chapter in this book asked a \emph{structural} question about a
tree you already had.  Chapter~29 (\chref{c31}) was the last: it taught you to
\emph{extract} a tree from a graph that was one.  This chapter turns the arrow
around.  The input is no longer a specific tree; the input is a \emph{property},
and the question is \textbf{how many labelled trees satisfy it}.  That flips
every tool: instead of one DFS you build a bijection, instead of a fold you
count, and instead of a numerical answer you return a residue modulo a prime.

The chapter is built on three pillars, each its own section.  The \term{Pr\"ufer
bijection} (\(\secref{c32-prufer}\)) is a length-\(n-2\) sequence of labels from
\([n]=\{1,\dots,n\}\) that encodes a labelled tree on \([n]\) with a two-way
greedy inverse; every enumeration question about labelled trees collapses to
\emph{count sequences}, which is elementary.  \term{Cayley's formula}
(\(\secref{c32-cayley}\)) is the immediate corollary \(n^{n-2}\), but the real
payoff is in its variants -- the degree-sequence multinomial and the
\emph{forest-completion} formula \(n^{k-2}\prod_i s_i\) that counts how many
ways you can add \(k-1\) edges to a \(k\)-component forest to make it a single
tree.  The \term{Matrix-Tree theorem} (\(\secref{c32-matrix}\)) handles the
case where the graph is not complete: the number of its spanning trees is a
cofactor of its Laplacian, so a modular Gaussian elimination is the whole
algorithm.  A fourth short section, \term{linearity of expectation}
(\(\secref{c32-expect}\)) on a rooted
tree, is the probabilistic twin of counting: split
an expectation into vertex indicators, and \(\mathbb{E}\) becomes a sum of
reciprocals of depths.

\textbf{You will learn to}
\begin{itemize}[nosep]
  \item encode and decode the Pr\"ufer sequence of a labelled tree, read the
        degree of every vertex off the code, and count degree-constrained
        trees via a multinomial (\(\secref{c32-prufer}\)); \hfill\lv{3}
  \item state and use Cayley \(n^{n-2}\), rooted \(n^{n-1}\), and the
        forest-completion formula \(n^{k-2}\prod_i s_i\); solve CF 2190D
        \emph{Prufer Vertex}, which asks for the second-to-last peeling
        survivor of every completion, modulo \(998244353\)
        (\(\secref{c32-cayley}\)); \hfill\lv{5}
  \item build the Laplacian, delete a row and column, and evaluate the
        determinant modulo a prime to count spanning trees of an arbitrary
        graph (\(\secref{c32-matrix}\)); \hfill\lv{4}
  \item decompose an expectation into per-vertex indicators, use
        \(\Pr[v\text{ is picked before its parent}]=1/\mathrm{depth}(v)\),
        and solve CF 280C \emph{Game on Tree} in one DFS
        (\(\secref{c32-expect}\)); \hfill\lv{4}
  \item tell when a counting problem is \emph{labelled} (this chapter) and
        when it is \emph{unlabelled} (Ch.28 \chref{c22}), and use the
        orbit-size bridge to reconcile the two (\(\secref{c32-boundary}\)).
        \hfill\lv{5}
\end{itemize}

\textbf{Prerequisites.} Chapter~13 (\chref{c10}) -- the bottom-up tree DP fold,
which the Summit uses to combine a polynomial across a rooted tree in one
recursion.  Chapter~12 (\chref{c09}) -- depth and path length in a rooted tree,
which \(\secref{c32-expect}\) uses verbatim.  Chapter~17 (\chref{c28}) -- DSU
component census, which \(\secref{c32-cayley}\) uses to read the \(k\) and
\(s_i\) of the input forest.  Chapter~29 (\chref{c31}) -- the \textbf{previous
chapter}; it ends at "here is the tree hiding inside the graph"; this one opens
at "here is how many trees match a description."
\end{chaptertheorybox}

% ---------------- Reading paths + ladder position ----------------
\begin{chaptertheorybox}{How to Read This Chapter (three paths)}
Level badges \(\lv3\)--\lv{6}: \(\lv3\) is hand-encoding a labelled tree to its
Pr\"ufer code by dropping smallest leaves; \(\lv4\) is coding the decode
greedy and running the Laplacian determinant on a small graph, and applying
linearity of expectation to CF 280C; \(\lv5\) is the forest-completion formula
and CF 2190D; \(\lv6\) is the Summit (CF 917D) -- count labelled trees sharing
exactly \(k\) edges with a given one, via the completion formula and a
tree-DP polynomial.

\textbf{Bijection path (\(\lv3\)--\lv{5}, ``count sequences'').}
\secref{c32-prufer} in full -- encode, decode, degree-from-code -- capped by
\secref{c32-cayley}'s Cayley and forest-completion formulas.  Stop here if
your problem reduces to a closed form on \([n]\).

\textbf{Linear-algebra path (\(\lv4\)--\lv{5}, ``count spanning trees of a
given graph'').} \secref{c32-matrix}: Laplacian, cofactor, modular determinant,
and the sanity checks against Cayley on \(K_n\) and against \(n\) on \(C_n\).
This is the branch you take when the graph is \emph{given} and you cannot
reduce to a Pr\"ufer-style sequence.

\textbf{Probability path (\(\lv4\)--\lv{6}, ``expected value on a tree'').}
\secref{c32-expect} -- indicator decomposition and the depth-reciprocal
formula, then solved CF 280C in one DFS.  Read \(\secref{c32-boundary}\) after
either branch to know when you are outside this chapter's method.
\end{chaptertheorybox}

\begin{figure}[H]
\centering
\begin{tikzpicture}[diagram, x=8.5mm, y=6.5mm]
  \foreach \st/\lab in {1/lv1,2/lv2,3/lv3,4/lv4,5/lv5,6/lv6} {
    \fill[\lab!22, draw=\lab, line width=0.8pt] (\st-1,0) rectangle (\st,\st*0.52);
    \node[font=\scriptsize\bfseries, text=\lab] at (\st-0.5,\st*0.52-0.42) {L\st};}
  \node[annot, anchor=south west, font=\scriptsize]
    at (0.05,0.62) {prereq: degree/depth (Ch.12), tree DP (Ch.13), DSU census (Ch.17)};
  \node[annot, anchor=south west, font=\scriptsize]
    at (2.05,1.66) {30.1 Prufer encode / decode by hand; degree from multiplicity};
  \node[annot, anchor=south west, font=\scriptsize]
    at (3.05,2.18) {30.3 Laplacian + modular det | 30.4 linearity, CF 280C};
  \node[annot, anchor=south west, font=\scriptsize, text=lv5]
    at (4.05,2.70) {30.2 Cayley + forest completion -- CF 2190D};
  \node[annot, anchor=south west, font=\scriptsize, text=lv6]
    at (5.05,3.22) {30.5 boundary; Summit CF 917D (exactly k shared edges)};
\end{tikzpicture}
\caption{Chapter~30 on the six-step ladder: hand-encoding the Pr\"ufer code
of the running tree is \(\lv3\); the coded decode, the Laplacian determinant,
and CF 280C's depth-reciprocal sum are \(\lv4\); CF 2190D's
per-vertex completion count is \(\lv5\); the Summit (CF 917D) is step~6.}
\end{figure}

% ---------------- The running example ----------------
\subsection*{The Running Example: A Labelled Tree on Six Vertices}

Every section up to the Summit runs on the same 6-vertex labelled tree \(T\),
fixed once here and authoritative (if a later trace disagrees, this figure
wins).  Vertices are \([6]=\{1,2,3,4,5,6\}\); edges, 1-based, are
\[
  1\!-\!5,\quad 2\!-\!5,\quad 5\!-\!3,\quad 3\!-\!4,\quad 4\!-\!6 .
\]
The Pr\"ufer code of \(T\) is \([5,5,3,4]\): drop the smallest leaf in turn
(\(1\), write its neighbour \(5\); then \(2\), write \(5\); then \(5\) itself
becomes a leaf after \(1,2\) are gone, but the smallest remaining leaf is
\(\min\{5,6\}=5\) if we hold on the definition -- \(\secref{c32-prufer}\) does
the exact bookkeeping), and the length-\(n-2=4\) sequence is \([5,5,3,4]\).
Reading degrees off the code: a vertex appears \((\deg-1)\) times, so
\(\deg(5)=3\), \(\deg(3)=\deg(4)=2\), \(\deg(1)=\deg(2)=\deg(6)=1\).  The
degree sum is \(3+2+2+1+1+1=10=2(n-1)\), as it must be.

\begin{figure}[H]
\centering
\begin{tikzpicture}[diagram, scale=0.95, transform shape,
   tnode/.style={circle, draw=primaryblue!70, fill=primaryblue!8, line width=0.8pt,
                 minimum size=8mm, font=\small\bfseries, inner sep=1pt}]
  \node[tnode] (v1) at (-2.6,1.2) {1};
  \node[tnode] (v2) at (-2.6,-1.2) {2};
  \node[tnode] (v5) at (0,0) {5};
  \node[tnode] (v3) at (2.4,0) {3};
  \node[tnode] (v4) at (4.4,-0.8) {4};
  \node[tnode] (v6) at (6.2,0.4) {6};
  \draw[primaryblue!80, line width=1.1pt] (v1)--(v5) (v2)--(v5) (v5)--(v3) (v3)--(v4) (v4)--(v6);
  \node[annot, font=\scriptsize, anchor=north, text width=9cm] at (1.2,-1.8)
    {leaves: \(1,2,6\) \quad deg 3: \(5\) \quad deg 2: \(3,4\)};
\end{tikzpicture}
\caption{The running tree \(T\).  Its Pr\"ufer code is \([5,5,3,4]\), and
\(\sum_v 1/\mathrm{depth}(v)\) from the root \(1\) is \(1/1+1/2+1/3+1/3+1/4+1/5
= 157/60\) (used in \(\secref{c32-expect}\)).  This object is fixed for the
whole chapter.}
\label{fig:c32-tree}
\end{figure}

\begin{figure}[H]
\centering
\begin{tikzpicture}[diagram, scale=1.0, transform shape,
  box/.style={rectangle, draw=primaryblue!70, fill=primaryblue!8,
              minimum width=9mm, minimum height=8mm,
              font=\small\bfseries, inner sep=2pt}]
  \node[box] (a1) at (0,0) {5};
  \node[box] (a2) at (1.2,0) {5};
  \node[box] (a3) at (2.4,0) {3};
  \node[box] (a4) at (3.6,0) {4};
  \node[annot, font=\scriptsize, anchor=east] at (-0.6,0) {code:};
  \node[annot, font=\scriptsize, anchor=west] at (4.4,0) {\((n-2)=4\) entries};
  \node[annot, font=\scriptsize, anchor=north] at (1.2,-0.7) {\(\deg(5)-1=2\) copies};
  \node[annot, font=\scriptsize, anchor=north] at (2.4,-0.7) {\(\deg(3)-1=1\)};
  \node[annot, font=\scriptsize, anchor=north] at (3.6,-0.7) {\(\deg(4)-1=1\)};
\end{tikzpicture}
\caption{The Pr\"ufer code \([5,5,3,4]\) laid out.  The number of times a
label appears in the code is exactly \(\deg-1\); every counting application in
this chapter is a way of exploiting that one fact.}
\label{fig:c32-code}
\end{figure}

Two traps return throughout.  The \textbf{labelled-vs-unlabelled trap}: this
chapter counts labelled trees, and two trees with the same shape but a
different vertex-numbering are different objects.  Chapter~28 (\chref{c22})
counts \emph{shapes}; on \(n=6\) there are exactly 6 unlabelled tree shapes
but \(6^4=1296\) labelled ones.  Reconciling the two is the orbit-size
argument in \(\secref{c32-boundary}\); confusing them is the most common wrong
answer on Cayley-style problems.  The \textbf{code-is-a-sequence trap} is
subtler: the Pr\"ufer code is a sequence \emph{with structure}, not a bag of
labels.  Decoding needs the greedy "smallest label not in the residual
suffix" step to be well-defined; encoding drops the smallest \emph{leaf},
which is a strictly stronger notion than smallest label.  Every wrong
implementation of Pr\"ufer decode gets these two confusions somewhere.

\subsection*{Which Job? Reading the Statement}

\begin{center}
\renewcommand{\arraystretch}{1.25}
\setlength{\tabcolsep}{3pt}
\begin{tabular}{@{}p{0.30\linewidth} p{0.20\linewidth} p{0.28\linewidth} >{\raggedright\arraybackslash}p{0.13\linewidth}@{}}
\toprule
\textbf{The statement says\dots} & \textbf{Tool} & \textbf{Why it works} & \textbf{Section}\\
\midrule
"how many trees on \([n]\) with property \(P\)"; \(P\) reads off the code & Pr\"ufer bijection & trees \(\leftrightarrow\) length-\((n-2)\) sequences & \(\secref{c32-prufer}\)\\
"count completions of a forest"; "rooted trees"; "with prescribed degrees" & Cayley variants & \(n^{n-2}\), \(n^{k-2}\prod s_i\), multinomial & \(\secref{c32-cayley}\)\\
"spanning trees of this graph"; graph given explicitly & Matrix-Tree & Laplacian cofactor = count & \(\secref{c32-matrix}\)\\
"expected value of a game on a tree" & Linearity of expectation & per-vertex indicators \(\to\) reciprocals of depths & \(\secref{c32-expect}\)\\
\bottomrule
\end{tabular}
\end{center}

\begin{insightbox}[The whole chapter in one sentence]
A labelled tree on \([n]\) is a sequence of \(n-2\) labels in disguise; every
enumerative question you can ask about such trees is either (a) a count of
those sequences with a filter, or (b) a linear-algebra determinant of a
Laplacian, and the entire craft of this chapter is picking which of the two
the problem actually is -- then applying linearity of expectation when the
answer is a real number, not an integer.
\end{insightbox}

% ---------------- Taxonomy figure ----------------
\begin{figure}[H]
\centering
\begin{tikzpicture}[diagram, scale=0.9, transform shape,
   mnode/.style={circle, draw=primaryblue!70, fill=primaryblue!8, line width=0.7pt,
                 minimum size=13mm, font=\scriptsize\bfseries, inner sep=1pt, align=center}]
  \node[mnode] (root) at (0,0) {count\\ trees};
  \node[mnode] (lab)  at (-4,-1.6) {labelled\\ on \([n]\)};
  \node[mnode] (sub)   at (1.5,-1.6) {of a given\\ graph};
  \node[mnode] (prf)  at (-5.6,-3.2) {Prufer\\ code};
  \node[mnode] (cay)   at (-2.6,-3.2) {Cayley /\\ completion};
  \node[mnode] (mt)   at (0.4,-3.2) {Matrix-\\ Tree};
  \node[mnode] (exp)  at (3.4,-3.2) {Linearity\\ of \(\mathbb{E}\)};
  \draw (root)--(lab)--(prf);
  \draw (lab)--(cay);
  \draw (root)--(sub)--(mt);
  \draw (sub)--(exp);
\end{tikzpicture}
\caption{Where the four methods sit. A \emph{labelled} enumeration on \([n]\)
is either a Pr\"ufer-code argument (\(\secref{c32-prufer}\)) or a Cayley
variant (\(\secref{c32-cayley}\)).  Counting inside a \emph{given graph} is
either Matrix-Tree (\(\secref{c32-matrix}\)) or, if the answer is
probabilistic, a linearity-of-expectation argument (\(\secref{c32-expect}\)).}
\label{fig:c32-taxonomy}
\end{figure}

\subsection*{Who Owns What}

\begin{center}
\renewcommand{\arraystretch}{1.2}
\setlength{\tabcolsep}{3pt}
\begin{tabular}{@{}p{0.28\linewidth} p{0.40\linewidth} p{0.24\linewidth}@{}}
\toprule
\textbf{Object} & \textbf{Owned here?} & \textbf{Elsewhere}\\
\midrule
Pr\"ufer code, bijection encode/decode & \textbf{Yes}, \(\secref{c32-prufer}\) -- with worked example & --\\
Cayley \(n^{n-2}\), forest completion \(n^{k-2}\prod s_i\) & \textbf{Yes}, \(\secref{c32-cayley}\) -- proved from the bijection & --\\
Matrix-Tree theorem, Laplacian cofactor & \textbf{Yes}, \(\secref{c32-matrix}\) & --\\
Linearity of expectation on rooted trees & \textbf{Yes}, \(\secref{c32-expect}\) & --\\
Depth / path length & no -- used as a black box & Ch.12 (\chref{c09})\\
Bottom-up tree DP fold & no -- the fold is reused & Ch.13 (\chref{c10})\\
DSU component census & no -- used once for forest reading & Ch.17 (\chref{c28})\\
Unlabelled tree shapes & no -- different question & Ch.28 (\chref{c22})\\
Minimum spanning tree & no -- a different "trees in a graph" & Ch.31 (\chref{c17})\\
\bottomrule
\end{tabular}
\end{center}

\subsection*{Recognition Guide}

\begin{recognitionbox}[Turning a counting question into a Pr\"ufer answer]
\reckey{``how many trees on \(n\) vertices satisfy \(\ldots\)'', ``count labelled
trees sharing exactly \(k\) edges with a given tree'', ``expected number of
random-deletion steps until empty'', ``given degrees, how many trees'', ``add
edges to make a forest connected''}
\recuse{
  \item \textbf{Property reads off degrees or the sequence itself} -- Pr\"ufer
        bijection (\(\secref{c32-prufer}\)) turns it into a multinomial; Cayley
        variants (\(\secref{c32-cayley}\)) finish it.
  \item \textbf{Property depends on the graph's edge set} -- Matrix-Tree
        (\(\secref{c32-matrix}\)): Laplacian cofactor.
  \item \textbf{Answer is a real number, not an integer} -- linearity of
        expectation (\(\secref{c32-expect}\)); the per-vertex indicator is
        almost always \(1/\)depth.
}
\recnotuse{
  \item \textbf{The question is about tree \emph{shapes} not labelings} --
        that is Ch.28 (\chref{c22}) tree isomorphism and AHU.
  \item \textbf{The question is about a specific tree, not a family} -- that
        is Ch.12/13/26/29, not this chapter.
}
One-line test: \emph{if the answer is a number mod a prime or a floating
value and the input is a \emph{description} rather than a graph, you are in
this chapter; build the bijection, apply the formula, or split the
expectation.}
\end{recognitionbox}

\subsection*{Constraint \(\to\) Technique}

\begin{center}
\renewcommand{\arraystretch}{1.2}
\setlength{\tabcolsep}{3pt}
\begin{tabular}{@{}p{0.34\linewidth} p{0.34\linewidth} >{\raggedright\arraybackslash}p{0.20\linewidth}@{}}
\toprule
\textbf{Constraint signature} & \textbf{Move} & \textbf{Section}\\
\midrule
\(n\le 10^6\), "how many trees with these degrees" & multinomial from code & \(\secref{c32-prufer}\)\\
forest on \(n\le 2\cdot 10^5\), \(k\) components & \(n^{k-2}\prod s_i\) mod \(p\) & \(\secref{c32-cayley}\)\\
graph on \(n\le 100\), count spanning trees & Laplacian det mod \(p\) & \(\secref{c32-matrix}\)\\
tree on \(n\le 10^5\), expected steps to empty & \(\sum_v 1/\mathrm{depth}(v)\) & \(\secref{c32-expect}\)\\
\(n\le 100\), count trees sharing exactly \(k\) edges & completion + DP + inversion & \(\probref{c32-summit}\)\\
\bottomrule
\end{tabular}
\end{center}

\subsection*{Notation Used Only Here}

\([n]=\{1,2,\dots,n\}\) is the label set.  \(\mathcal{T}_n\) is the set of
labelled trees on \([n]\); \(|\mathcal{T}_n|=n^{n-2}\) by Cayley.  \(\mathrm{code}(T)\) is the
Pr\"ufer code of \(T\) (a sequence of \(n-2\) labels from \([n]\));
\(\mathrm{decode}(c)\) is its inverse.  \(\deg_T(v)\) is the degree of \(v\) in
\(T\); \(m_v(T)=\deg_T(v)-1\) is the \emph{multiplicity} of \(v\) in
\(\mathrm{code}(T)\).  For a graph \(G\), \(L(G)=D-A\) is its Laplacian, and
\(L^{(i)}(G)\) is \(L(G)\) with row and column \(i\) deleted.  For a rooted
tree, \(\mathrm{depth}(v)\) counts the vertices on the root-\(v\) path
(inclusive of both endpoints), so \(\mathrm{depth}(\text{root})=1\).  The
Summit uses \(f(k)\) for the count of completions whose shared-edge set with
the given tree contains a fixed \(k\)-subset, and \(g(k)\) for the count that
shares \emph{exactly} \(k\) edges -- these two are related by a binomial
inversion.
